\section*{Chapter 1: Introduction}
\setcounter{section}{1} \setcounter{subsection}{0}
\phantomsection \label{chap:introduction} \addcontentsline{toc}{section}{\textit{Chapter 1: Introduction}}

Recent advances in Large Language Models (LLMs) have enabled instruction following, reasoning, and interaction with external environments, forming the foundation for autonomous agentic systems \cite{yao2023react, liu2024agentbench, schick2023toolformer}.
LLM-based agents are now deployed across a wide range of applications, from direct-acting assistants to sophisticated multi-step reasoning and problem-solving systems \cite{park2023generative, qin2024toolllm, mei2024aios, liu2026agentos}.

\subsection{Problem Statement}

Tool-using, stateful AI agents expose a fundamental mismatch between modern agent frameworks and the abstractions provided by conventional operating systems \cite{packer2023memgpt, mei2024aios}.
While conventional operating systems effectively manage processes, files, and network resources, MCP-style workflows are organized around higher-level concepts such as semantic tool descriptions, session context, approval boundaries, and cross-step state.
This mismatch leads to fragile control planes when enforcement remains entirely in user space. 
Frameworks such as OpenClaw~\cite{openclaw2026} demonstrate that granting agents execution capability and persistent context can lead to security risks including prompt injection, privilege escalation, and cross-step exploitation. 

A single gateway process is typically responsible for interpreting tool manifests, binding requests to identities, and authorizing execution, effectively placing the control plane in user space. 
Such a design introduces a single point of failure and expands the trusted computing base, making it vulnerable to crashes, logic bugs, or compromised control paths~\cite{peter2014arrakis}.
Two properties become hard to enforce: whether a tool's declared behaviour matches what actually runs, and whether approval state survives a gateway crash~\cite{qin2024toolllm}.
Identity binding under replay and session confusion~\cite{greshake2023prompt} is a third, related problem.

This leads to a concrete question: which control-plane invariants in MCP-style tool invocation benefit from kernel-level enforcement, and can that enforcement improve security and resilience without making the system unacceptably slow?

\subsection{Our Solution}

\texttt{linux-mcp} addresses this by moving admission control, approval state, binding checks, and durable observability into a Linux kernel module, while leaving semantic parsing and tool execution in user space.
This is not a full AI-native operating system (AI-native OS); the kernel role is deliberately narrow, enforcing only the invariants user space cannot reliably hold.

The architecture has four layers: an LLM client (\texttt{llm-app}), a manifest-aware gateway (\texttt{mcpd}), a Generic-Netlink-based kernel control plane (\texttt{kernel\_mcp}), and resident tool services over Unix domain sockets.
Every tool invocation follows a fixed path through all four layers; no request reaches a tool service without first obtaining an explicit kernel admission decision.
The component responsibilities are detailed in \S\ref{sec:system-architecture}, and the rationale for the kernel--user-space split is developed in \S\ref{sec:design-alternatives}.

This design does not reduce the trusted footprint of \texttt{mcpd}: the gateway remains fully trusted and is responsible for populating the kernel registries.
What it provides instead is a \emph{partitioning} of responsibility: registry entries, approval tickets, and binding checks reside in the kernel, protected by \texttt{CAP\_NET\_ADMIN}-gated Generic Netlink access control, and remain durable across daemon failures and inaccessible to unprivileged processes.
Gateway logic bugs therefore no longer threaten the critical state transitions; the kernel enforces them independently of the daemon's correctness.
Semantic parsing and tool execution remain in user space, where flexibility matters most.

Rather than replacing existing Linux security mechanisms, \texttt{linux-mcp} proposes a kernel-assisted control plane that complements user-space flexibility.
AIOS-style proposals offer strong guarantees at the cost of replacing OS infrastructure; user-space frameworks offer broad deployment at the cost of the structural weaknesses identified in Chapter~\ref{chap:background}.
The position here is narrower: move only the invariants that require protection from the gateway itself into the kernel, and leave everything else in user space.
The central hypothesis is that this targeted placement improves binding safety and failure resilience without imposing prohibitive performance costs; Chapter~\ref{chap:evaluation} tests this claim empirically.

\paragraph{Comparison with user-space systems.}
To contextualise the design, we compare \texttt{linux-mcp} with two representative user-space tool-invocation frameworks: OpenClaw~\cite{openclaw2026}, which uses a monolithic user-space gateway, and OpenHarness~\cite{openharness2026}, which integrates control logic into the agent process itself.
Table~\ref{tab:comparison} highlights the key architectural differences.

\begin{table}[t]
\centering
\small
\setlength{\tabcolsep}{6pt}
\begin{tabularx}{\textwidth}{l X X X}
\toprule
\textbf{Feature} & \textbf{linux-mcp} & \textbf{OpenClaw} & \textbf{OpenHarness} \\
\midrule

Control-plane placement
& Kernel-assisted
& User-space gateway
& In-process (same PID as agent) \\[4pt]

Trust boundary
& Split (user\,+\,kernel)
& Monolithic (user-space)
& Python process boundary \\[4pt]

Approval enforcement
& Ticket-based (kernel)
& User-space logic
& In-memory checker \\[4pt]

Binding correctness
& Kernel-verified (UDS peer creds)
& Gateway-managed
& No OS-level binding \\[4pt]

Failure resilience
& State persists in kernel
& Lost on gateway failure
& Lost on process restart \\[6pt]

\midrule
\multicolumn{4}{c}{\textit{Security model}} \\
\midrule

\multicolumn{4}{p{\dimexpr\textwidth-2\tabcolsep\relax}}{\textbf{linux-mcp:} TCB partitioned; kernel holds durable control-plane state, inaccessible to unprivileged processes.} \\[4pt]

\multicolumn{4}{p{\dimexpr\textwidth-2\tabcolsep\relax}}{\textbf{OpenClaw:} Monolithic TCB; all control-plane state in user-space process.} \\[4pt]

\multicolumn{4}{p{\dimexpr\textwidth-2\tabcolsep\relax}}{\textbf{OpenHarness:} No OS-level enforcement boundary; sensitive-path deny lists enforced only within the Python interpreter.} \\

\bottomrule
\end{tabularx}
\caption{Key design differences between \texttt{linux-mcp}, OpenClaw, and OpenHarness.}
\label{tab:comparison}
\end{table}

\noindent



\subsection{Main Contributions}

This paper makes the following contributions:
\begin{itemize}

    \item We design and implement \texttt{kernel\_mcp}, a Linux kernel module that maintains tool and agent registries, enforces per-invocation admit/deny/defer decisions, and manages approval-ticket lifecycle under \texttt{CAP\_NET\_ADMIN}-gated Generic Netlink access.
    This makes control-plane state durable across daemon failure and unreachable by unprivileged processes, without placing any JSON parsing or tool execution in the kernel.

    \item We design and implement \texttt{mcpd}, a manifest-aware user-space gateway that loads tool manifests, derives semantic hash identifiers, reconciles tool and agent state with the kernel registry, binds sessions to OS-level UDS peer credentials, validates request payloads against JSON schemas, and forwards tool invocations through the kernel mediation path.
    \texttt{mcpd} is the sole component that bridges manifest semantics and runtime enforcement, keeping this complexity out of both the kernel and the LLM client.

    \item We conduct a controlled functional evaluation against pure and hardened user-space baselines across latency, throughput, attack resistance, and failure scenarios.
    The evaluation demonstrates that \texttt{linux-mcp} enforces the specified control-plane invariants against a targeted attack set ($n = 30$ runs), quantifies the latency cost of kernel placement, and shows that kernel-held approval state survives daemon failure in a way neither user-space baseline can match.
\end{itemize}
